\section{A compatible signature bound}\label{sec:signature}

The discriminant relations must hold simultaneously on the actual
tree. A square discriminant at one height, without a compatible
choice of signs at all vertices, would not prove the required bound.

\begin{proposition}[Signature upper bound]\label{prop:upper}
Over the original field $k$, the inverse-image tree admits a
compatible labelling for which $G_n\subseteq E_n$ for every
$n\geq1$.
\end{proposition}

\begin{proof}
For a vertex $v$ with value $x_v$, order its three children and let
$\Delta_v$ be their Vandermonde product. By
\eqref{eq:discproduct},
\[
 \Delta_v^2=a^3x_v,\qquad
 \left(\prod_{w\text{ child of }v}\Delta_w\right)^2=a^9x_v.
\]
All these quantities are nonzero because the generic roots are
distinct and no vertex value is zero. Consequently
\[
 \frac{\prod_{w\text{ child of }v}\Delta_w}{a^3\Delta_v}
 \in\{1,-1\}.
\]
Choose the order of the first-level children. Having chosen all
orders through one height, choose orders one height lower so that
\begin{equation}\label{eq:compatible-vandermonde}
 \prod_{w\text{ child of }v}\Delta_w=a^3\Delta_v
\end{equation}
at every newly treated parent $v$. If the sign is initially wrong,
an odd permutation of the children of one child $w$ changes exactly
one factor on the left. The correction is independent for different
parents. It does not alter a Vandermonde already fixed higher in
the tree, so it preserves the previously imposed relations.
Induction produces one labelling of the infinite tree.

For $g\in G_n$, let $s_v(g)$ be the sign of its permutation of the
three ordered children of $v$ onto those of $g(v)$. Apply $g$ to
\eqref{eq:compatible-vandermonde} and compare it with that identity
at $g(v)$. Since $a$ is fixed, the comparison gives
\[
 s_v(g)=\prod_{w\text{ child of }v}s_w(g).
\]
These are exactly the local defining relations of $E_n$ at all
vertices with two remaining levels. They prove $G_n\subseteq E_n$.
The construction took place over $k$ itself, without adjoining
constants.
\end{proof}

We record two elementary consequences of the classical group
definition. The order formula is
\cite[Proposition~2.2]{benedetto2017cubic}; its short recurrence is
included to make the later degree comparison explicit.

\begin{lemma}[Order and bottom actions]\label{lem:order}
The groups in \eqref{eq:Edefinition} satisfy
\begin{equation}\label{eq:order-recurrence}
 |E_1|=6,\qquad |E_{n+1}|=3|E_n|^3,
 \qquad
 \frac{|E_{n+1}|}{|E_n|}
   =2^{2\cdot3^{n-1}}3^{3^n}.
\end{equation}
For every bottom parent of $T_n$, $E_n$ contains the subgroup
$\alt_3$ permuting that parent's three children and fixing every
other vertex.
\end{lemma}

\begin{proof}
The product-sign homomorphism on $E_n^3\rtimes\sym_3$ is onto:
a root transposition with identity sections has value $-1$.
Its kernel therefore has order $6|E_n|^3/2=3|E_n|^3$.
Solving this recurrence gives
\[
 |E_n|=2^{3^{n-1}}3^{(3^n-1)/2},
\]
and division gives the final expression in
\eqref{eq:order-recurrence}. An independent bottom three-cycle
has positive local sign everywhere, so it satisfies every local
product-sign relation. The same holds for its square and the
identity.
\end{proof}

The bottom actions in Lemma~\ref{lem:order} will be used only after
the induction has proved $G_n=E_n$ at the current height. They will
not be assumed for $G_{n+1}$.

